\section*{Capítulo \ref{ch:b3:heterocycles} --- Química heterocíclica}

\begin{solution}{exo:b3:heterocycles:1}
Oxazol: el O aporta dos (un par solitario en el sistema $\pi$ y otro en el plano), el N uno, y los
tres carbonos tres: 6. Tiazol: lo mismo con el S: 6. Pirimidina: uno de cada uno
de los seis átomos, con los pares solitarios de los dos nitrógenos en el plano: 6. Purina: nueve átomos;
el nitrógeno N--H aporta dos y los demás uno cada uno: 10, un número de Hückel
($4n + 2$, $n = 2$).
\end{solution}

\begin{solution}{exo:b3:heterocycles:2}
Piperidina (11.1) $>$ DMAP (9.6) $>$ imidazol (7.0) $>$ piridina (5.2) $\gg$ pirrol (alrededor de $-4$,
protonado en un carbono).
\end{solution}

\begin{solution}{exo:b3:heterocycles:3}
2-Bromopirrol (el pirrol es tan reactivo que se polibroma con facilidad);
2-bromofurano; 3-bromoindol.
\end{solution}

\begin{solution}{exo:b3:heterocycles:4}
2,5-Dimetilfurano; 1,2,5-trimetilpirrol; 2,5-dimetiltiofeno.
\end{solution}

\begin{solution}{exo:b3:heterocycles:5}
2-Metoxipiridina. La adición del metóxido en C2 sitúa la carga negativa del
\omterm{def:b3:heterocycles:snar}{complejo de Meisenheimer} en el nitrógeno; en C3 solo puede repartirse sobre carbonos, de modo que el
intermedio, y el estado de transición que lo precede, tienen una energía mucho más alta.
\end{solution}

\begin{solution}{exo:b3:heterocycles:6}
2-Aminopiridina (en forma de su sal sódica, que se protona después en el tratamiento). El ion amiduro
se adiciona en C2, lo que da un aducto $\sigma$ aniónico con H y $\ce{NH2}$ en C2 y la carga
negativa en el nitrógeno del anillo; este pierde un hidruro, que toma un protón del
grupo amino y se va como $\ce{H2}$.
\end{solution}

\begin{solution}{exo:b3:heterocycles:7}
El oxígeno del N-óxido devuelve un par solitario al anillo: las estructuras
resonantes llevan carga negativa en C2, C4 y C6. Los electrófilos atacan C4 (C2
está cerca del nitrógeno con carga positiva). Después se elimina el oxígeno con
tricloruro de fósforo, lo que da la 4-nitropiridina.
\end{solution}

\begin{solution}{exo:b3:heterocycles:8}
2-Hidroxipiridina (un OH en C2) y 2-piridona (N--H y C=O). La piridona conserva seis
electrones $\pi$ en el anillo: su estructura resonante zwitteriónica, con $\mathrm{O^-}$ y un
$\mathrm{N^+}$ que comparte su par solitario, es un piridinio-olato; es una amida cuyo
par solitario del nitrógeno completa el sexteto, y los fuertes enlaces de hidrógeno de C=O y N--H
la favorecen en agua.
\end{solution}

\begin{solution}{exo:b3:heterocycles:9}
2-Nitrobenzaldehído, dos equivalentes de acetoacetato de metilo y amoníaco: el
bloqueador de los canales de calcio nifedipino.
\end{solution}

\begin{solution}{exo:b3:heterocycles:10}
La enhidrazina orientada hacia el grupo $\ce{CH2}$ (interna, más sustituida) da el
2,3-dimetilindol; la orientada hacia el grupo metilo (terminal) da el 2-etilindol.
Los ácidos débiles dan sobre todo el primero, a través de la enhidrazina más estable;
los ácidos más fuertes aumentan la proporción de 2-etilindol.
\end{solution}

\begin{solution}{exo:b3:heterocycles:11}
En la $\mathrm S_\mathrm NAr$, la etapa determinante de la velocidad es la adición, que no rompe el enlace C--X;
el flúor, más electronegativo, estabiliza más el estado de transición aniónico
y el \omterm{def:b3:heterocycles:snar}{complejo de Meisenheimer}. En la $\mathrm S_\mathrm N2$, el enlace C--X se rompe en el
estado de transición, y el enlace C--I, débil, es el que se rompe más deprisa.
\end{solution}

\begin{solution}{exo:b3:heterocycles:12}
Piridina: tres señales en la razón 2\,:\,1\,:\,2, la primera hacia \qty{8.6}{ppm}. Pirrol: dos señales iguales
a 6.7 y \qty{6.2}{ppm} y un N--H ancho hacia \qty{8.1}{ppm}. Furano: dos señales iguales a
7.4 y \qty{6.4}{ppm}, sin N--H. Tiofeno: dos señales iguales a 7.33 y \qty{7.12}{ppm}, próximas
entre sí, entre las del furano.
\end{solution}

\begin{solution}{pb:b3:heterocycles:1}
\textbf{1.} $\ce{C6H5NHNH2 + (CH3)2CO -> C6H5NHN=C(CH3)2 + H2O}$.

\textbf{2.} El nitrógeno de la amina se adiciona al carbono carbonílico y después se pierde agua, como en
la formación de una imina; el ácido activa el grupo carbonilo y convierte el OH del aducto
en un grupo saliente.

\textbf{3.} El $\ce{NH2}$ terminal: el par solitario del otro nitrógeno está conjugado con el
anillo fenilo y más impedido.

\textbf{4.} $\qty{108.0}{g/mol}$; $\qty{10.0}{g}/\qty{108.0}{g/mol} = \qty{0.0926}{mol}$.

\textbf{5.} La hidrazona $\ce{C9H12N2}$, \qty{148.0}{g/mol}: $\qty{0.0926}{mol} \times \qty{148.0}{g/mol} = \qty{13.7}{g}$.

\textbf{6.} $\ce{C6H5-NH-NH-C(CH3)=CH2}$.

\textbf{7.} El carbono \textit{orto} del anillo, el carbono \textit{ipso}, los dos nitrógenos, el carbono de la hidrazona
y el $\ce{CH2}$ terminal.

\textbf{8.} Se rompe el enlace N--N; se forma un enlace C--C entre el carbono \textit{orto} y el $\ce{CH2}$.

\textbf{9.} Seis electrones, con todos los componentes \omterm{def:b3:pericyclic:faciality}{suprafaciales}: permitida por vía térmica.

\textbf{10.} Como ácido de Lewis, cataliza la tautomerización, se une a un nitrógeno y
debilita el enlace N--N, y ayuda a que salga el amoníaco.

\textbf{11.} El carbono \textit{orto} se ha vuelto $sp^3$, con un hidrógeno: la pérdida de ese protón
restablece el anillo bencénico, una fuerza impulsora importante.

\textbf{12.} El $\ce{NH2}$ anilínico formado ataca el carbono de la imina y cierra el
anillo de cinco miembros (una 2-aminoindolina).

\textbf{13.} El amoníaco; la impulsa la formación del anillo de pirrol aromático del indol.

\textbf{14.} $\ce{C9H12N2 -> C9H9N + NH3}$.

\textbf{15.} \qty{131.0}{g/mol}.

\textbf{16.} $\qty{0.0926}{mol} \times \qty{131.0}{g/mol} = \qty{12.1}{g}$.

\textbf{17.} En C3, la posición más nucleófila del indol, libre aquí; el ataque en ella
mantiene intacto el anillo bencénico en el intermedio.

\textbf{18.} $\ce{C6H5NHNH-C(=CH2)-CH2CH3}$ (terminal) y $\ce{C6H5NHNH-C(CH3)=CH-CH3}$ (interna).

\textbf{19.} 2-Etilindol y 2,3-dimetilindol.

\textbf{20.} La interna está más sustituida; los ácidos fuertes aumentan la proporción de la
terminal.

\textbf{21.} El 2,3-dimetilindol, a través de la enhidrazina más estable y más
sustituida.

\textbf{22.} El acetaldehído se condensa consigo mismo, y la vía de la enhidrazina da rendimientos
bajos; el indol se prepara en su lugar a partir de la hidrazona del ácido pirúvico, que da el
ácido indol-2-carboxílico, que después se descarboxila.

\textbf{23.} Hay que separar una mezcla de regioisómeros, lo que cuesta rendimiento; se prefiere una
cetona simétrica o una vía que fije la regioquímica.

\textbf{24.} 10.0 g de fenilhidrazina pueden dar como máximo \qty{12.1}{g} de 2-metilindol.
\end{solution}
